% concatenated from main.tex
\documentclass{amsart}
\usepackage{amsmath,amssymb,amsthm, amsfonts,color}

\newtheorem{theorem}{Theorem}
\newtheorem{lemma}[theorem]{Lemma}
\newtheorem{proposition}[theorem]{Proposition}
\theoremstyle{definition}
\newtheorem{definition}[theorem]{Definition}
\theoremstyle{remark}
\newtheorem{remark}[theorem]{Remark}

\makeatletter
\newcommand\thankssymb[1]{\textsuperscript{\@fnsymbol{#1}}}
\makeatother

\newcommand{\Fp}{\mathbb{F}_p}
\newcommand{\Zpos}{\mathbb{Z}^{+}}
\DeclareMathOperator{\cc}{cc}
\DeclareMathOperator{\rank}{rank}
\DeclareMathOperator{\GL}{GL}

\title[List recovery for random low-rate linear codes]{List recovery for random low-rate linear codes}
\date{}

\begin{document}

\author[Isaac M Hair \and Amit Sahai]{Isaac M Hair\thankssymb{1} \and Amit Sahai\thankssymb{2}}

\thanks{\thankssymb{1} UCSB, UCLA. \texttt{isaacmhair@gmail.com}}
\thanks{\thankssymb{2} UCLA. \texttt{sahai@cs.ucla.edu}}



\begin{abstract}
We prove a list recovery guarantee for random low-rate linear codes over sufficiently large prime fields.  For fixed dimension \(d\), error fraction \(\alpha\), and accuracy parameter \(\varepsilon\), a random \(d\)-dimensional linear code \(C \subseteq \Fp^n\) is, with high probability, \((\alpha,\ell,\frac{1+\varepsilon}{1-\alpha}\ell)\)-list recoverable simultaneously for all input list sizes \(\ell\le 2^{O_{\alpha, \varepsilon, d}(n/\log n)}\). The proof is inspired by a work of Matou\v{s}ek, P\v{r}\'{\i}v\v{e}tiv\'{y}, and \v{S}kovro\v{n} on reconstructing point sets from their projections. It combines a deterministic graph-theoretic certificate, a nonvanishing determinant criterion, and the Schwartz--Zippel lemma.  We also give a lower bound showing that any linear code \(C \subseteq \Fp^n\) of dimension at least two \emph{cannot} be \((\alpha,\ell,\frac{1+\varepsilon}{1-\alpha}\ell)\)-list recoverable for feasible list sizes $\ell \geq 2^{\Omega_{\alpha, \varepsilon}(n)}$. In this sense, our result is nearly optimal.
\end{abstract}

\maketitle




\section{Introduction}

In this work, we study list recovery of random linear codes over large prime fields. We use the following standard formulation~\cite{ReschVenkitesh}.

\begin{definition}[$(\alpha, \ell, L)$-List Recoverable]\label{def:list-recoverable}
    A code $C \subseteq \Sigma^n$ is said to be $(\alpha, \ell, L)$-list recoverable if for all $S_1, \ldots, S_n \subseteq \Sigma$ with $\vert S_1 \vert = \ldots = \vert S_n\vert = \ell$, there are at most $L$ codewords $x \in C$ such that
    \[\vert \{ i : x_i \not\in S_i\}\vert \leq \alpha n.\]
\end{definition}

The case $\ell=1$ is ordinary list decoding. List recovery has its roots in soft-decision and list-decoding work for concatenated and Reed--Solomon-type codes, including early concatenated-code bounds and interpolation-based algorithms \cite{ZyablovPinsker,GuruswamiSudan}. It is now a standard primitive in coding theory, for example in expander-code, tensor-code, and local-list-recovery constructions \cite{HemenwayRonZewiWootters}; see also the recent survey \cite{ReschVenkitesh}. Random linear codes form a particularly important test case. Their list-decoding behavior was studied by Guruswami--H\aa stad--Kopparty and Wootters \cite{GuruswamiHastadKopparty,Wootters}. For list recovery, Rudra and Wootters gave average-radius tools that apply in low-rate large-field regimes \cite{RudraWootters}. More recently, Guruswami--Li--Mosheiff--Resch--Silas--Wootters established sharp list-size phenomena and separations from fully random codes near capacity \cite{GLMRSW}, and other works have considered specific linear families, small-field random linear codes, zero-rate thresholds, and Brascamp--Lieb-type combinatorial bounds \cite{LiShagrithaya,DoronMosheiffReschRibeiro,ReschYuanZhang,BrakensiekChenDharZhang}. The present paper studies a complementary, essentially zero-rate random-linear-code regime in which the dimension is fixed while the block length tends to infinity.

\textbf{Our Results.}
In this paper, we study list recovery for random linear codes of fixed dimension $d$ and growing block length $n$. We show that for all constants $\alpha>0$ and $\epsilon>0$, over sufficiently large prime fields, a uniformly random linear code $C \subseteq \Fp^n$ of dimension $d$ is list recoverable with near-optimal output list size $L=\frac{1+\varepsilon}{1-\alpha}\ell$ simultaneously for every input list size $\ell\le 2^{\delta n/\log n}$ for some constant $\delta>0$. The lower bound following the main theorem shows that this dependence on $n$ is close to the best possible: for every linear code of dimension at least two, the same conclusion cannot generally be extended to all
input list sizes of order $2^{\Theta(n)}$. (We remark that the lower bound follows from standard techniques in the coding-theory literature.)

\begin{theorem}[Main theorem]\label{thm:main}
    For all $\alpha \in (0,1)$, $\varepsilon \in (0,1)$, and $d \in \mathbb{Z}^+$, there exist $\delta > 0$, $n_0 \in \mathbb{Z}^+$, where 
$n_0=O_{\alpha,\varepsilon}(d\log d)$ for large enough $d$, and a computable function $f : \mathbb{Z}^+ \rightarrow \mathbb{Z}^+$ such that the following holds. Let $n \geq n_0$ be any integer, and let $p \geq f(n)$ be any prime. With probability at least $1 - \varepsilon$, a random linear code $C \subseteq \Fp^n$ of dimension $d$ is $(\alpha, \ell, \frac{1+\varepsilon}{1-\alpha}\ell)$-list recoverable for all positive integers $\ell \leq 2^{\delta n/\log n}$.
\end{theorem}

\begin{theorem}[Exponential lower bound]\label{thm:lower}
    For all $\alpha \in (0,1)$ and $\varepsilon \in (0,1)$, there exist $\delta > 0$, $n_0 \in \mathbb{Z}^+$, and a computable function $f : \mathbb{Z}^+ \rightarrow \mathbb{Z}^+$ such that the following holds. Let $n \geq n_0$ be any integer, and let $p \geq f(n)$ be any prime. Let $C \subseteq \Fp^n$ be any linear code of dimension at least two. For every feasible list size $\ell$, i.e. every integer $\ell$ with $2^{\delta n}\le \ell\le p$, $C$ is not $(\alpha, \ell, \frac{1+\varepsilon}{1-\alpha}\ell)$-list recoverable.
\end{theorem}

\textbf{Statement on AI use.}
The human authors are fully responsible for the contents of this paper.
The questions posed and answered in this work are fully due to the human authors.
The human authors obtained a less optimal form of the results of this paper entirely using human reasoning. However, we then posed our question -- without providing any information about our approach -- to the UCLA Moonshot AI for Math harness~\cite{MoonshotHarness} which made extensive use of queries to GPT 5.5Pro. 
The proof obtained by the harness was superior to our human-devised one both quantitatively and \emph{qualitatively}. We also note that GPT5.5Pro on its own was not able to prove this theorem when the authors attempted to query it directly. Almost all of this paper was directly written by the harness after some post-processing by GPT 5.5Pro.

\textbf{Proof Overview.}
We start by explaining the connection between list recovery and certain colored graph problems. Consider a linear code $C=\operatorname{im}M\subseteq\Fp^n$ that is not list recoverable with the desired parameters, where $M:\Fp^d\to\Fp^n$ has coordinate forms $\lambda_1,\ldots,\lambda_n$. Then there is a set $A\subseteq\Fp^d$ of more than $\frac{1+\varepsilon}{1-\alpha}\ell$ messages whose codewords land in prescribed lists $S_1,\ldots,S_n$ in all but an $\alpha$-fraction of the coordinates. For a fixed coordinate $i$, the points of $A$ that land in $S_i$ are distributed among at most $\ell$ fibers of $\lambda_i$. Since $|A|$ is larger than the target output list size, a counting argument shows that many coordinates force many pairs of points of $A$ to lie in common fibers. Pairing such points inside the fibers gives a matching, and we color each edge by the coordinate $i$ that produced it. Thus a failure of list recovery yields an edge-colored graph on the vertex set $A$: colors are coordinates, color classes are matchings, and an edge $\{a,a'\}$ of color $i$ records the linear equation $\lambda_i(a-a')=0$.

Once the problem has been converted into this graph language, we use a graph-theoretic and linear-algebraic strategy adapted from a work of Matou\v{s}ek, P\v{r}\'{\i}v\v{e}tiv\'y, and \v{S}kovro\v{n} \cite{MPS}. Section~2 proves the graph-theoretic ingredient: if many colors each form a large matching, then on some vertex subset one can find $d$ spanning trees whose color sets are pairwise disjoint. The proof first shows that a random set of colors connects a graph under a suitable color-expansion hypothesis, and then obtains that expansion from a maximal-density induced subgraph.

Section~3 turns these color-disjoint spanning trees into an algebraic certificate. The tree equations form a square linear system in the point differences, and the disjointness of the colors gives a determinant polynomial that is not identically zero. Whenever this determinant evaluates nonzero, the equations force all points in the certificate to coincide. A union bound over certificates, together with the Schwartz--Zippel lemma, shows that random coordinate forms avoid all small determinant certificates with high probability. Section~4 records the deterministic consequence: if the coordinate forms avoid these certificates, then no large set of messages can lie close to the prescribed coordinate lists. Section~5 applies this to a random linear map $\Fp^d\to\Fp^n$ and then conditions on full rank to obtain a uniformly random $d$-dimensional subspace. Finally, Section~6 proves the lower bound by passing to a two-dimensional subcode and building a generalized arithmetic box whose coordinate projections all have size at most $\ell$ while the box itself contains more than $\frac{1+\varepsilon}{1-\alpha}\ell$ codewords.

\section{Preliminaries and colored graphs}

For a positive integer \(m\), write $[m]=\{1,\ldots,m\}$. 
All graphs are finite. Multigraphs are allowed unless explicitly excluded. A graph is loopless if it has no loops. In an edge-colored multigraph, a color class is a matching if no two edges of that color share a vertex. For an edge-colored graph \(H\) with color set \([m]\) and \(I\subseteq[m]\), let \(H[I]\) denote the spanning subgraph containing exactly the edges whose colors lie in \(I\). For a graph \(J\), let \(\cc(J)\) be the number of connected components, isolated vertices included.

For \(b\ge 2\), define
\[
\Lambda(b):=\max\{1,\log b\,\log\log(4b)\}.
\]
The function \(\Lambda\) is nondecreasing on \([2,\infty)\): indeed, if
\[
g(b)=\log b\,\log\log(4b),
\]
then
\[
g'(b)=\frac{1}{b}\log\log(4b)+\frac{\log b}{b\log(4b)}>0
\qquad(b\ge2),
\]
and the maximum of two nondecreasing functions is nondecreasing.

We use the following standard form of the Schwartz--Zippel lemma without proof.

\begin{lemma}[Schwartz--Zippel \cite{Schwartz,Zippel}]\label{lem:schwartz-zippel}
Let \(F\) be a field, let \(S\subseteq F\) be finite and nonempty, and let
\[
P\in F[x_1,\ldots,x_N]
\]
be a nonzero polynomial of total degree at most \(D\). If \(u_1,\ldots,u_N\) are independent and uniformly distributed in \(S\), then
\[
\Pr(P(u_1,\ldots,u_N)=0)\le \frac{D}{|S|}.
\]
\end{lemma}

\begin{lemma}[Random colors give connectivity]\label{lem:random-colors}
Let \(H\) be a loopless edge-colored multigraph on \(w\ge2\) vertices with color set \([m]\). Suppose that, for some \(\gamma>0\), every nonempty \(A\subseteq V(H)\) with \(|A|\le w/2\) has at least
\[
\gamma m\log\frac{w}{|A|}
\]
distinct colors crossing the cut \((A,V(H)\setminus A)\). Let \(0<\eta<1\), and let \(I\) be a uniformly random \(T\)-element subset of \([m]\). If \(T\le m\) and
\[
T\ge \frac4\gamma\left(\log\log(4w)+\log(1/\eta)+1\right),
\]
then
\[
\Pr(H[I]\text{ is connected})\ge 1-\eta.
\]
\end{lemma}

\begin{proof}
The singleton cuts imply \(\gamma\log w\le1\), so \(0<\gamma/4<1\). Expose the \(T\) colors one at a time without replacement. Let \(I_j\) be the first \(j\) colors, and set
\[
X_j=\cc(H[I_j]),\qquad Z_j=\log X_j.
\]
We prove that
\[
\mathbb E[Z_{j+1}\mid I_j]\le \left(1-\frac\gamma4\right)Z_j
\]
whenever \(j<T\).

Condition on \(I_j\) and put \(x=X_j\). If \(x=1\), there is nothing to prove, so assume \(x\ge2\). Let \(C_1,\ldots,C_x\) be the current components, and call \(C_i\) small if \(|C_i|\le w/2\). No exposed color crosses a current component. Hence, for a small component \(C\), the next color crosses \(C\) with conditional probability at least
\[
\frac{\gamma m\log(w/|C|)}{m-j}\ge \gamma\log\frac{w}{|C|}.
\]
If \(Y\) is the number of small components crossed by the next color, then
\[
\mathbb E[Y\mid I_j]\ge \gamma\sum_{C\text{ small}}\log\frac{w}{|C|}.
\]
We claim that
\[
\sum_{C\text{ small}}\log\frac{w}{|C|}\ge \frac{x}{2}\log x. \tag{1}
\]
If all components are small, Jensen's inequality gives
\[
\sum_{i=1}^x\log\frac{w}{|C_i|}\ge x\log x.
\]
Otherwise there is exactly one component larger than \(w/2\). If the remaining \(r=x-1\) components have total size \(s\le w/2\), Jensen's inequality gives
\[
\sum_{C\text{ small}}\log\frac{w}{|C|}
\ge r\log\frac{rw}{s}
\ge r\log(2r)
\ge \frac{x}{2}\log x,
\]
since \(r\ge x/2\) and \(2r\ge x\). This proves (1).

Let \(\Delta_j=X_j-X_{j+1}\) be the decrease in the number of components after adding the next color. In the auxiliary graph whose vertices are the current components and whose edges are the newly added crossing edges, if \(t\) vertices are nonisolated, then the number of components drops by at least \(t/2\). Since every component counted by \(Y\) is nonisolated, \(\Delta_j\ge Y/2\). Therefore
\[
\mathbb E[\Delta_j\mid I_j]
\ge \frac{\gamma x\log x}{4}.
\]
Using
\[
\log\frac{x}{x-\Delta_j}\ge \frac{\Delta_j}{x},
\]
we get
\[
\mathbb E[Z_j-Z_{j+1}\mid I_j]
\ge \frac{\gamma}{4}\log x
=\frac\gamma4 Z_j.
\]
This is the desired one-step inequality.

Iterating gives
\[
\mathbb E Z_T
\le \left(1-\frac\gamma4\right)^T\log w
\le e^{-\gamma T/4}\log w.
\]
If \(H[I]\) is disconnected, then \(Z_T\ge\log2\), so Markov's inequality and the assumed lower bound on \(T\) yield
\[
\Pr(H[I]\text{ is disconnected})
\le \frac{\log w}{\log2}\,e^{-\gamma T/4}
\le \eta\frac{\log w}{e(\log2)\log(4w)}
\le \eta.
\]
\end{proof}

\begin{lemma}[Large color matchings force disjoint trees]\label{lem:many-matchings}
Fix \(\beta>0\) and \(d\in\Zpos\). There is a constant \(C_{\rm gr}=C_{\rm gr}(\beta,d)\) such that the following holds for every integer \(b\ge2\).

Let \(G\) be a loopless edge-colored multigraph on \(b\) vertices with color set \([m]\). Suppose every color class is a matching and has at least \(\beta b\) edges. If
\[
m\ge C_{\rm gr}\Lambda(b),
\]
then there is a vertex set \(W\subseteq V(G)\), \(|W|\ge2\), and there are \(d\) spanning trees
\[
T_1,\ldots,T_d\subseteq G[W]
\]
whose color sets are pairwise disjoint.
\end{lemma}

\begin{proof}
For integers \(s\ge2\), set
\[
T(s)=
\left\lceil
\frac{4\log s}{\beta}
\left(\log\log(4s)+\log(2d)+1\right)
\right\rceil.
\]
Since \(s\ge2\), \(d\ge1\), and \(\Lambda(s)\ge1\),
\[
\log s\log\log(4s)\le \Lambda(s),
\qquad
\log s\le O(\Lambda(s)),
\]
with an absolute implicit constant. Therefore
\[
T(s)\le O_\beta(\log(2d)\Lambda(s))
\qquad\text{for all }s\ge2.
\]
Choose \(C_{\rm gr}\ge d\), with \(C_{\rm gr}=O_\beta(d\log(2d))\), so that
\[
dT(s)\le C_{\rm gr}\Lambda(s)
\qquad\text{for all }s\ge2.
\]

Let \(G\) be as in the statement. For \(U\subseteq V(G)\), let \(e(U)\) be the number of edges of \(G[U]\), counted with multiplicity. Since each color class has at least \(\beta b\) edges,
\[
e(V(G))\ge \beta mb.
\]
Choose \(W\subseteq V(G)\), \(w:=|W|\ge2\), maximizing
\[
\rho(U)=\frac{e(U)}{|U|\log |U|}
\]
over all \(U\) with \(|U|\ge2\), and put \(\rho=\rho(W)\). Then
\[
\rho\ge \frac{e(V(G))}{b\log b}\ge \frac{\beta m}{\log b}.
\]
Let \(H=G[W]\). We verify the expansion hypothesis of Lemma~\ref{lem:random-colors}. Take a nonempty \(A\subseteq W\) with \(a:=|A|\le w/2\), and let \(D=W\setminus A\). By maximality of \(W\), with the interpretation \(e(U)=0=\rho |U|\log |U|\) when \(|U|=1\),
\[
e(A)\le \rho a\log a,
\qquad
e(D)\le \rho |D|\log |D|.
\]
Thus the number \(e(A,D)\) of edges crossing the cut satisfies
\[
\begin{aligned}
e(A,D)
&\ge \rho\bigl(w\log w-a\log a-(w-a)\log(w-a)\bigr)\\
&=\rho\left(a\log\frac wa+(w-a)\log\frac{w}{w-a}\right)\\
&\ge \rho a\log\frac wa.
\end{aligned}
\]
A single color is a matching, so it contributes at most \(a\) edges to this cut. Hence at least
\[
\rho\log\frac wa
\ge \frac{\beta m}{\log b}\log\frac wa
\]
distinct colors cross the cut. Lemma~\ref{lem:random-colors} applies to \(H\) with
\[
\gamma=\frac{\beta}{\log b}.
\]

Since
\[
m\ge C_{\rm gr}\Lambda(b)\ge dT(b),
\]
choose pairwise disjoint random sets \(I_1,\ldots,I_d\subseteq[m]\) with
\[
|I_1|=\cdots=|I_d|=T(b),
\]
for instance by taking consecutive blocks in a random permutation of \([m]\). Each \(I_r\) is marginally uniform. Because \(w\le b\), the definition of \(T(b)\) and Lemma~\ref{lem:random-colors} with \(\eta=1/(2d)\) give
\[
\Pr(H[I_r]\text{ is connected})\ge1-\frac1{2d}
\]
for each \(r\). By the union bound, with positive probability all \(H[I_r]\) are connected. Fix such a choice and select a spanning tree \(T_r\) of each \(H[I_r]\). The color sets of these trees are pairwise disjoint because the sets \(I_r\) are.
\end{proof}

\section{Algebraic certificates}

\begin{lemma}[Disjoint trees give a nonzero determinant]\label{lem:determinant}
Let \(F\) be a field. Let \(w\ge2\) and \(d\ge1\), and let \(G\) be an edge-colored multigraph on vertex set \([w]\), with colors in \([m]\). For each color \(i\in[m]\), introduce formal coefficients
\[
a_{i,1},\ldots,a_{i,d}
\]
and the formal linear form
\[
\lambda_i(z_1,\ldots,z_d)=\sum_{j=1}^d a_{i,j}z_j.
\]
Suppose \(G\) contains \(d\) spanning trees \(T_1,\ldots,T_d\) whose color sets are pairwise disjoint. Orient every edge of \(T_1\cup\cdots\cup T_d\) arbitrarily, counting edges with multiplicity. Let \(R\) be the \(d(w-1)\times d(w-1)\) matrix whose rows are indexed by these oriented edges \(e=(u,v)\), whose columns are indexed by \((s,j)\in[w-1]\times[d]\), and whose entries are
\[
R_{e,(s,j)}
=
(\mathbf 1_{s=u}-\mathbf 1_{s=v})a_{\chi(e),j},
\]
where \(\chi(e)\) is the color of \(e\). Then \(\det R\) is a nonzero polynomial in the variables \(a_{i,j}\).

Consequently, after any specialization \(a_{i,j}\mapsto \alpha_{i,j}\in F\) for which \(\det R(\alpha)\ne0\), if points \(q_1,\ldots,q_w\in F^d\) satisfy
\[
\sum_{j=1}^d \alpha_{\chi(e),j}(q_u-q_v)_j=0
\]
for every oriented edge \(e=(u,v)\) of \(T_1\cup\cdots\cup T_d\), then
\[
q_1=\cdots=q_w.
\]
\end{lemma}

\begin{proof}
Let \(C_r\) be the set of colors appearing in \(T_r\). The sets \(C_1,\ldots,C_d\) are pairwise disjoint. Specialize the variables by setting \(a_{i,r}=1\) and \(a_{i,j}=0\) for \(j\ne r\) whenever \(i\in C_r\), and setting all unused-color variables to \(0\). With rows grouped by the trees and columns grouped by the coordinate \(j\), the matrix becomes block diagonal. The \(r\)th block is, up to row signs and row and column permutations, the reduced incidence matrix of the spanning tree \(T_r\) with the column for vertex \(w\) deleted.

A reduced incidence matrix of a tree has determinant \(\pm1\): root the tree at \(w\), orient edges toward the root, order non-root vertices by decreasing distance from \(w\), and order rows by the corresponding child vertices. The resulting reduced incidence matrix is triangular with diagonal entries \(1\). Changing orientations or orders only changes the sign. Therefore, under the specialization above,
\[
\det R=\pm1,
\]
so \(\det R\) is not the zero polynomial over \(F\).

For the consequence, suppose \(\det R(\alpha)\ne0\). Let
\[
Q=\bigl((q_s-q_w)_j\bigr)_{s\in[w-1],\,j\in[d]}\in F^{d(w-1)}.
\]
The edge equations are exactly
\[
R(\alpha)Q=0,
\]
because the terms involving \(q_w\) cancel and there is no column for the base vertex \(w\). Since \(R(\alpha)\) is invertible, \(Q=0\), and hence every \(q_s\) equals \(q_w\).
\end{proof}

\begin{definition}[Tree certificates and good linear forms]\label{def:good}
Fix integers \(B\ge2\), \(m\ge1\), and \(d\ge1\), and a field \(F\). A tree certificate up to size \(B\) with \(m\) colors and \(d\) trees is a tuple
\[
\mathcal T=(w,T_1,\ldots,T_d)
\]
such that \(2\le w\le B\), each \(T_r\) is a colored spanning tree on vertex set \([w]\) with colors in \([m]\), and the color sets of \(T_1,\ldots,T_d\) are pairwise disjoint. 
A colored spanning tree may use the same color on more than one edge; its color set is the set of colors that appear at least once.
For each certificate, fix an arbitrary orientation and ordering of the rows and columns in Lemma~\ref{lem:determinant}, and let \(P_{\mathcal T}\) be the resulting determinant polynomial.

Linear forms
\[
\lambda_i(z_1,\ldots,z_d)=\sum_{j=1}^d \alpha_{i,j}z_j,
\qquad i\in[m],
\]
are called good up to \(B\) if
\[
P_{\mathcal T}(\alpha)\ne0
\]
for every tree certificate \(\mathcal T\) up to size \(B\).
\end{definition}

\begin{proposition}[Random forms avoid all small certificates]\label{prop:random-good}
Let \(B\ge2\), \(m\ge1\), \(d\ge1\), and let \(p\) be a prime. Choose \(m\) random linear forms
\[
\lambda_i(z_1,\ldots,z_d)=\sum_{j=1}^d a_{i,j}z_j,
\qquad i\in[m],
\]
with all coefficients \(a_{i,j}\) independent and uniformly distributed in \(\Fp\). Then
\[
\Pr(\lambda_1,\ldots,\lambda_m\text{ are good up to }B)
\ge
1-\frac{dB^2(B^2m)^{dB}}{p}.
\]
\end{proposition}

\begin{proof}
For fixed \(w\), a colored spanning tree on \([w]\) is overcounted by an ordered list of \(w-1\) triples consisting of two endpoints and a color, giving at most \((B^2m)^{w-1}\) choices. Thus the number of \(d\)-tuples of colored spanning trees for this \(w\) is at most
\[
(B^2m)^{d(w-1)}\le(B^2m)^{dB}.
\]
There are at most \(B\) choices for \(w\), and the disjointness condition only decreases the count, so the number of certificates is at most
\[
B(B^2m)^{dB}.
\]
For a fixed certificate \(\mathcal T=(w,T_1,\ldots,T_d)\), Lemma~\ref{lem:determinant} gives a nonzero determinant polynomial \(P_{\mathcal T}\) over \(\Fp\). Its matrix has size \(d(w-1)\) and entries of degree at most \(1\), so
\[
\deg P_{\mathcal T}\le d(w-1)\le dB.
\]
By Lemma~\ref{lem:schwartz-zippel}, this polynomial vanishes at the random coefficient array with probability at most \(dB/p\). The union bound over all certificates gives the stated estimate.
\end{proof}


\section{The deterministic implication}

\begin{lemma}[Good forms rule out large approximate list intersections]\label{lem:deterministic}
Let $\alpha,\varepsilon\in(0,1)$ and $d\in\Zpos$, and put
\[
K=\frac{1+\varepsilon}{1-\alpha},
\qquad
\mu=\frac{(1-\alpha)\varepsilon}{2(1+\varepsilon)},
\qquad
\beta=\frac\mu2,
\qquad
\theta=\frac{\mu}{1-\mu}.
\]
Let $C_{\rm gr}=C_{\rm gr}(\beta,d)$ be the constant from Lemma~\ref{lem:many-matchings}. Let $B\ge2$, let $n\in\Zpos$ satisfy
\[
n\ge \frac{C_{\rm gr}}{\theta}\Lambda(B),
\]
let $p$ be a prime, and suppose that linear forms
\[
\lambda_1,\ldots,\lambda_n:\Fp^d\to\Fp
\]
are good up to $B$. Then there do not exist a set $A\subseteq\Fp^d$, a positive integer $\ell$, and sets
\[
S_1,\ldots,S_n\subseteq\Fp,
\qquad |S_1|=\cdots=|S_n|=\ell,
\]
such that
\[
2\le |A|\le B,
\qquad
|A|>K\ell,
\qquad
\left|\{i:\lambda_i(a)\notin S_i\}\right|\le \alpha n
\quad\text{for every }a\in A.
\]
\end{lemma}

\begin{proof}
Assume such $A$, $\ell$, and $S_1,\ldots,S_n$ exist, and put $b=|A|$. For each coordinate $i$, let
\[
A_i=\{a\in A:\lambda_i(a)\in S_i\},
\qquad
 a_i=|A_i|.
\]
The hypothesis on the number of bad coordinates gives
\[
\sum_{i=1}^n a_i\ge (1-\alpha)nb.
\]
For each $i$, partition $A_i$ into the nonempty fibers of $\lambda_i$. Since $\lambda_i(A_i)\subseteq S_i$, there are at most $\ell$ such fibers. Pair points arbitrarily inside each fiber, leaving at most one point unpaired per fiber. This gives a matching $M_i$ on the vertex set $A$ such that
\[
|M_i|\ge \frac12(a_i-\ell)_+,
\qquad (x)_+:=\max\{x,0\}.
\]
Therefore
\[
\sum_{i=1}^n |M_i|
\ge \frac12\sum_{i=1}^n(a_i-\ell)_+
\ge \frac12\left(\sum_{i=1}^n a_i-n\ell\right)
\ge \frac n2((1-\alpha)b-\ell).
\]
Since $b>K\ell=\frac{1+\varepsilon}{1-\alpha}\ell$, we have
\[
(1-\alpha)b-\ell
>
\frac{(1-\alpha)\varepsilon}{1+\varepsilon}b,
\]
and hence
\[
\sum_{i=1}^n |M_i|>\mu nb.
\]
Let
\[
I=\{i\in[n]: |M_i|\ge \beta b\}.
\]
Each matching has size at most $b/2$, so
\[
\mu nb
<\sum_{i=1}^n |M_i|
\le |I|\frac b2+(n-|I|)\beta b.
\]
Using $\beta=\mu/2$ and cancelling $b$ gives
\[
|I|>\frac{\mu}{1-\mu}n=\theta n.
\]
In particular,
\[
|I|\ge C_{\rm gr}\Lambda(B).
\]

Form a loopless edge-colored multigraph $G$ on vertex set $A$ by inserting the matching $M_i$ in color $i$ for every $i\in I$. Each color class is a matching of size at least $\beta b$. Since $b\le B$ and $\Lambda$ is nondecreasing, Lemma~\ref{lem:many-matchings}, applied after relabeling the colors in $I$, gives a subset $W\subseteq A$, $w:=|W|\ge2$, and $d$ spanning trees
\[
T_1,\ldots,T_d\subseteq G[W]
\]
whose color sets are pairwise disjoint.

Relabel $W$ as $[w]$, and write the corresponding points of $\Fp^d$ as
\[
q_1,\ldots,q_w.
\]
If an edge of color $i$ is oriented as $(u,v)$, then its endpoints were paired inside a single fiber of $\lambda_i$, so
\[
\lambda_i(q_u-q_v)=0.
\]
The trees $T_1,\ldots,T_d$ form a tree certificate up to size $B$ using colors from $[n]$. Since the forms are good up to $B$, the associated determinant is nonzero at their coefficient values. Lemma~\ref{lem:determinant} implies
\[
q_1=\cdots=q_w,
\]
contradicting the fact that $W$ consists of at least two distinct points of the set $A$.
\end{proof}

\section{Proof of the main theorem}

\begin{proof}[Proof of Theorem~\ref{thm:main}]
Fix $\alpha,\varepsilon\in(0,1)$ and $d\in\Zpos$. Put
\[
K=\frac{1+\varepsilon}{1-\alpha},
\qquad
\mu=\frac{(1-\alpha)\varepsilon}{2(1+\varepsilon)},
\qquad
\beta=\frac\mu2,
\qquad
\theta=\frac{\mu}{1-\mu},
\]
and let $C_{\rm gr}=C_{\rm gr}(\beta,d)$. By Lemma~\ref{lem:many-matchings}, we may take
\[
C_{\rm gr}=O_{\alpha,\varepsilon}(d\log(2d)).
\]
Choose $0<c=c(\alpha,\varepsilon)\le1$ sufficiently small, and set
\[
\delta=\frac{c}{C_{\rm gr}}.
\]
We choose $n_0\ge d+1$ with $n_0=O_{\alpha,\varepsilon}(C_{\rm gr})$ so that, for every $N\ge n_0$,
\[
N\ge \frac{C_{\rm gr}}{\theta}\Lambda\left(\left\lceil K2^{\delta N/\log N}\right\rceil+1\right).
\]
Indeed, for $x\ge0$ there is a constant $C_0=C_0(\alpha,\varepsilon)$ such that
\[
\Lambda\left(\left\lceil K2^x\right\rceil+1\right)
\le C_0(1+x)\log(e+x).
\]
Increasing $C_0$ if necessary, for all sufficiently large $N$ and all $0\le x\le N/\log N$ this gives
\[
\Lambda\left(\left\lceil K2^x\right\rceil+1\right)
\le C_0(1+x\log N).
\]
Since $\delta\le1$, we may apply this with $x=\delta N/\log N$. Taking $c\le \theta/(2C_0)$ and then taking $n_0\ge (2C_0/\theta)C_{\rm gr}$, enlarged if necessary to handle the fixed ``sufficiently large'' threshold, gives for all $N\ge n_0$
\[
\frac{C_{\rm gr}}{\theta}\Lambda\left(\left\lceil K2^{\delta N/\log N}\right\rceil+1\right)
\le \frac{C_0}{\theta}C_{\rm gr}+\frac{C_0c}{\theta}N
\le N.
\]
Thus, since $C_{\rm gr}=O_{\alpha,\varepsilon}(d\log(2d))$, we have $n_0=O_{\alpha,\varepsilon}(d\log d)$ for large enough $d$.

For $N\ge n_0$, define
\[
B_N=\left\lceil K2^{\delta N/\log N}\right\rceil+1
\]
and
\[
f(N)=\left\lceil \frac{2dB_N^2(B_N^2N)^{dB_N}}{\varepsilon}\right\rceil.
\]
Define $f$ arbitrarily on the finitely many positive integers $N<n_0$. This gives a computable function $f:\Zpos\to\Zpos$.

Now fix $n\ge n_0$ and a prime $p\ge f(n)$. Let $M$ be a random $n\times d$ matrix over $\Fp$, with all entries independent and uniform. Its rows define linear forms
\[
\lambda_1,\ldots,\lambda_n:\Fp^d\to\Fp,
\qquad
(Mt)_i=\lambda_i(t).
\]
Let $\mathcal R$ be the event $\rank M=d$. For a fixed nonzero $v\in\Fp^d$, the vector $Mv$ is uniformly distributed in $\Fp^n$, so
\[
\Pr(Mv=0)=p^{-n}.
\]
Thus
\[
\Pr(\mathcal R^\complement)
\le (p^d-1)p^{-n}
\le p^{d-n}
\le p^{-1},
\]
and in particular $\Pr(\mathcal R)\ge 1/2$.

Let $E$ be the event that the $n$ forms $\lambda_1,\ldots,\lambda_n$ are good up to $B_n$. Proposition~\ref{prop:random-good}, with $m=n$ and $B=B_n$, gives
\[
\Pr(E^\complement)
\le
\frac{dB_n^2(B_n^2n)^{dB_n}}{p}
\le \frac{\varepsilon}{2}.
\]
Consequently,
\[
\Pr(E^\complement\mid\mathcal R)
\le \frac{\Pr(E^\complement)}{\Pr(\mathcal R)}
\le \varepsilon.
\]

We claim that on $E\cap\mathcal R$, the code
\[
C_M=\operatorname{im}M\subseteq\Fp^n
\]
is $(\alpha,\ell,K\ell)$-list recoverable for every positive integer $\ell\le 2^{\delta n/\log n}$. Suppose not. Then for some such $\ell$ and some sets
\[
S_1,\ldots,S_n\subseteq\Fp,
\qquad |S_1|=\cdots=|S_n|=\ell,
\]
there are more than $K\ell$ codewords $x\in C_M$ such that
\[
\left|\{i:x_i\notin S_i\}\right|\le\alpha n.
\]
Choose
\[
b=\lfloor K\ell\rfloor+1
\]
distinct such codewords. Since $\mathcal R$ holds, the map $t\mapsto Mt$ is injective, so these codewords have a unique preimage set $A\subseteq\Fp^d$ of size $b$. Moreover,
\[
2\le b\le K\ell+1\le K2^{\delta n/\log n}+1\le B_n.
\]
For every $a\in A$, the corresponding codeword $Ma$ has at most $\alpha n$ coordinates outside the lists $S_i$, equivalently
\[
\left|\{i:\lambda_i(a)\notin S_i\}\right|\le\alpha n.
\]
This contradicts Lemma~\ref{lem:deterministic}, because $E$ says that $\lambda_1,\ldots,\lambda_n$ are good up to $B_n$ and the choice of $\delta,n_0$ gives
\[
n\ge \frac{C_{\rm gr}}{\theta}\Lambda(B_n).
\]
Thus $C_M$ satisfies the desired list-recovery conclusion on $E\cap\mathcal R$.

It remains to pass from random matrices to uniform random subspaces. Conditioned on $\mathcal R$, the columns of $M$ form an ordered basis of $\operatorname{im}M$. Conversely, every fixed $d$-dimensional subspace of $\Fp^n$ has exactly $|\GL_d(\Fp)|$ ordered bases, so $\operatorname{im}M$ conditioned on $\mathcal R$ is uniformly distributed over all $d$-dimensional subspaces of $\Fp^n$. Therefore the failure probability for a uniformly random $d$-dimensional linear code is at most
\[
\Pr(E^\complement\mid\mathcal R)\le\varepsilon,
\]
and the theorem follows.
\end{proof}

\section{Proof of the lower bound}

\begin{proof}[Proof of Theorem~\ref{thm:lower}]
Fix $\alpha,\varepsilon\in(0,1)$ and put
\[
K=\frac{1+\varepsilon}{1-\alpha}.
\]
Choose an integer $T>K+2$, set $\delta=\log_2 T$, and take $n_0=2$. Define, for example,
\[
f(N)=\left\lceil \max\{T^N,2^{2N+2}(K+1)\}\right\rceil.
\]
This is a computable function $\Zpos\to\Zpos$.

Let $n\ge n_0$, let $p\ge f(n)$ be prime, and let $C\subseteq\Fp^n$ be a linear code of dimension at least two. Fix an integer $\ell$ with
\[
T^n=2^{\delta n}\le \ell\le p.
\]
It suffices to construct lists $S_1,\ldots,S_n\subseteq\Fp$ of size $\ell$ and more than $K\ell$ codewords of $C$ that meet every list in every coordinate.

Choose a two-dimensional subcode $D\subseteq C$ and identify it with $\Fp^2$. Under this identification, the $i$th coordinate map on $D$ is a linear form
\[
\lambda_i:\Fp^2\to\Fp.
\]
The nonzero forms among the $\lambda_i$ span the dual of $\Fp^2$: otherwise some nonzero vector of $D$ would vanish in every coordinate, contradicting the fact that $D\subseteq\Fp^n$ is a two-dimensional subcode. Hence their kernels include at least two distinct one-dimensional subspaces. Let
\[
U_1,\ldots,U_r
\]
be the distinct kernels that occur among the nonzero coordinate forms, where $2\le r\le n$, and choose a nonzero vector $u_j\in U_j$ for each $j$.

We next choose side lengths. Set
\[
t_2=\cdots=t_r=T,
\qquad
t_1=\left\lfloor \frac{K\ell}{T^{r-1}}\right\rfloor+1,
\qquad
P=\prod_{j=1}^r t_j.
\]
Since $\ell\ge T^n\ge T^{r-1}$, we have
\[
P>K\ell.
\]
Also,
\[
\frac{P}{t_1}=T^{r-1}\le \ell,
\]
and for $j\ge2$,
\[
\frac{P}{t_j}=t_1T^{r-2}
\le \left(\frac{K\ell}{T^{r-1}}+1\right)T^{r-2}
=\frac{K}{T}\ell+T^{r-2}
\le \ell,
\]
because $T>K+2$ and $\ell\ge T^{r-1}$. Finally,
\[
P\le K\ell+T^{r-1}\le (K+1)\ell.
\]

We claim that the vectors $u_1,\ldots,u_r$ can be rescaled so that all sums
\[
\sum_{j=1}^r a_j s_j u_j,
\qquad
0\le a_j<t_j,
\]
are distinct. Indeed, for a nonzero difference vector
\[
\Delta=(\Delta_1,
\ldots,
\Delta_r),
\qquad
-(t_j-1)\le \Delta_j\le t_j-1,
\]
consider the bad equation
\[
\sum_{j=1}^r \Delta_j s_j u_j=0
\]
in the variables $s_1,
\ldots,s_r$. First note that all side lengths $t_j$ are less than $p$: this is clear for $j\ge2$, and for $j=1$ we use $r\ge2$ and $\ell\le p$ to get
\[
t_1\le \frac{K\ell}{T^{r-1}}+1\le \frac{Kp}{T}+1<p,
\]
where the final inequality follows from $T>K+2$ and $p\ge T^n\ge T^2$. Hence a nonzero integer $\Delta_j$ in the displayed range remains nonzero in $\Fp$. If $\Delta$ is supported on a single index, the bad equation has no solution with all $s_j\ne0$. If $\Delta$ has support of size at least two, then the involved vectors include two nonparallel vectors, so the equation imposes two independent linear conditions and has at most $p^{r-2}$ solutions in $\Fp^r$. The number of possible nonzero $\Delta$ is at most
\[
\prod_{j=1}^r (2t_j-1)
\le 2^rP
\le 2^n(K+1)\ell
\le 2^n(K+1)p.
\]
Thus the total number of bad choices of $(s_1,
\ldots,s_r)$ is at most
\[
2^n(K+1)p^{r-1}.
\]
Since $p\ge 2^{2n+2}(K+1)$ and $r\le n$, we have $p>2^{n+1}(K+1)$ and $p\ge2r$. Therefore
\[
2^n(K+1)p^{r-1}<\frac{p^r}{2}\le (p-1)^r,
\]
where the last inequality follows from $(1-1/p)^r\ge 1-r/p\ge 1/2$. Thus the bad choices do not cover all $(p-1)^r$ choices of nonzero scalars, so there is a choice of nonzero scalars $s_1,
\ldots,s_r$ for which all the displayed sums are distinct.

Fix such scalars and define
\[
A=\left\{\sum_{j=1}^r a_js_ju_j:0\le a_j<t_j\right\}\subseteq\Fp^2.
\]
Then $|A|=P>K\ell$. Under the identification $D\cong\Fp^2$, view $A$ as a set of codewords in the subcode $D\subseteq C$. If $\lambda_i=0$, then $|\lambda_i(A)|=1$. Otherwise, $\ker\lambda_i=U_j$ for some $j$, and the coefficient $a_j$ does not affect $\lambda_i$; hence
\[
|\lambda_i(A)|\le \prod_{h\ne j}t_h=\frac{P}{t_j}\le\ell.
\]
For each coordinate $i$, choose an $\ell$-element set $S_i\subseteq\Fp$ containing $\lambda_i(A)$, which is possible because $\ell\le p$. Every codeword in $A$ then lies in $S_i$ in every coordinate $i$. Therefore more than $K\ell$ codewords of $C$ have zero bad coordinates with respect to the lists $S_1,
\ldots,S_n$, so $C$ is not $(\alpha,\ell,K\ell)$-list recoverable.
\end{proof}

\section*{Acknowledgements}
This research was supported in part from a Simons Investigator Award, DARPA expMath award, Laude Moonshot grant, NSF grant 2333935, BSF grant 2022370, a Xerox Faculty Research Award, a Google Faculty Research Award, an Okawa Foundation Research Grant, and the Symantec Chair of Computer Science. This material is based upon work supported by the Defense Advanced Research Projects Agency through Award HR001126CE054.

\begin{thebibliography}{ZHC+26}

\bibitem[BCDZ25]{BrakensiekChenDharZhang}
Joshua Brakensiek, Yeyuan Chen, Manik Dhar, and Zihan Zhang,
\emph{Combinatorial Bounds for List Recovery via Discrete Brascamp--Lieb Inequalities},
arXiv:2510.13775, 2025.

\bibitem[DMRR25]{DoronMosheiffReschRibeiro}
Dean Doron, Jonathan Mosheiff, Nicolas Resch, and Jo\~ao Ribeiro,
\emph{List-Recovery of Random Linear Codes over Small Fields},
in \emph{Approximation, Randomization, and Combinatorial Optimization. Algorithms and Techniques (APPROX/RANDOM 2025)},
LIPIcs, vol.~353, 57:1--57:18, 2025.
doi:10.4230/LIPIcs.APPROX/RANDOM.2025.57.

\bibitem[GLM+22]{GLMRSW}
Venkatesan Guruswami, Ray Li, Jonathan Mosheiff, Nicolas Resch, Shashwat Silas, and Mary Wootters,
\emph{Bounds for List-Decoding and List-Recovery of Random Linear Codes},
IEEE Transactions on Information Theory \textbf{68} (2022), no.~2, 923--939.
doi:10.1109/TIT.2021.3127126.

\bibitem[GHK11]{GuruswamiHastadKopparty}
Venkatesan Guruswami, Johan H\aa stad, and Swastik Kopparty,
\emph{On the List-Decodability of Random Linear Codes},
IEEE Transactions on Information Theory \textbf{57} (2011), no.~2, 718--725.
doi:10.1109/TIT.2010.2095170.

\bibitem[GS99]{GuruswamiSudan}
Venkatesan Guruswami and Madhu Sudan,
\emph{Improved Decoding of Reed--Solomon and Algebraic-Geometry Codes},
IEEE Transactions on Information Theory \textbf{45} (1999), no.~6, 1757--1767.
doi:10.1109/18.782097.

\bibitem[HRW20]{HemenwayRonZewiWootters}
Brett Hemenway, Noga Ron-Zewi, and Mary Wootters,
\emph{Local List Recovery of High-Rate Tensor Codes and Applications},
SIAM Journal on Computing \textbf{49} (2020), no.~4, FOCS17-157--FOCS17-195.
doi:10.1137/17M116149X.

\bibitem[LS25]{LiShagrithaya}
Ray Li and Nikhil Shagrithaya,
\emph{Near-Optimal List-Recovery of Linear Code Families},
in \emph{Approximation, Randomization, and Combinatorial Optimization. Algorithms and Techniques (APPROX/RANDOM 2025)},
LIPIcs, vol.~353, 53:1--53:14, 2025.
doi:10.4230/LIPIcs.APPROX/RANDOM.2025.53.

\bibitem[MPS08]{MPS}
Ji\v{r}\'{\i} Matou\v{s}ek, Ale\v{s} P\v{r}\'{\i}v\v{e}tiv\'y, and Petr \v{S}kovro\v{n},
\emph{How Many Points Can Be Reconstructed from \(k\) Projections?},
SIAM Journal on Discrete Mathematics \textbf{22} (2008), no.~4, 1605--1623.
doi:10.1137/080715706.

\bibitem[RV25]{ReschVenkitesh}
Nicolas Resch and S. Venkitesh,
\emph{List Recoverable Codes: The Good, the Bad, and the Unknown (hopefully not Ugly)},
arXiv:2510.07597, 2025.

\bibitem[RYZ23]{ReschYuanZhang}
Nicolas Resch, Chen Yuan, and Yihan Zhang,
\emph{Zero-Rate Thresholds and New Capacity Bounds for List-Decoding and List-Recovery},
in \emph{50th International Colloquium on Automata, Languages, and Programming (ICALP 2023)},
LIPIcs, vol.~261, 99:1--99:18, 2023.
doi:10.4230/LIPIcs.ICALP.2023.99.

\bibitem[RW18]{RudraWootters}
Atri Rudra and Mary Wootters,
\emph{Average-Radius List-Recoverability of Random Linear Codes},
in \emph{Proceedings of the Twenty-Ninth Annual ACM-SIAM Symposium on Discrete Algorithms (SODA)},
2018, 644--662.
doi:10.1137/1.9781611975031.42.

\bibitem[Woo13]{Wootters}
Mary Wootters,
\emph{On the List Decodability of Random Linear Codes with Large Error Rates},
in \emph{Proceedings of the Forty-Fifth Annual ACM Symposium on Theory of Computing (STOC)},
2013, 853--860.
doi:10.1145/2488608.2488716.

\bibitem[ZP81]{ZyablovPinsker}
V. V. Zyablov and M. S. Pinsker,
\emph{List concatenated decoding},
Problemy Peredachi Informatsii \textbf{17} (1981), no.~4, 29--33.

\bibitem[Sch80]{Schwartz}
J. T. Schwartz,
\emph{Fast probabilistic algorithms for verification of polynomial identities},
Journal of the ACM \textbf{27} (1980), no.~4, 701--717.

\bibitem[Zip79]{Zippel}
R. Zippel,
\emph{Probabilistic algorithms for sparse polynomials},
in \emph{Symbolic and Algebraic Computation}, EUROSAM 1979,
Lecture Notes in Computer Science, vol.~72, Springer, 1979, 216--226.

\bibitem[ZHC+26]{MoonshotHarness}
Junyi Zhang\textsuperscript{*}, Xinjie He\textsuperscript{*}, Hyunsik Chae,
Ethan Ji, Eric Jiang, Rushil Raghavan, Yiwen Kou, Alex Taylor,
Kai-Wei Chang\textsuperscript{\textdagger}, Raghu Meka\textsuperscript{\textdagger},
Violet Peng\textsuperscript{\textdagger}, Amit Sahai\textsuperscript{\textdagger},
Terence Tao\textsuperscript{\textdagger}, and Wei Wang\textsuperscript{\textdagger}.
\newblock \emph{UCLA Moonshot Harness}.
\newblock 2026.
\newblock \textsuperscript{*}Co-first authors with equal contribution.
The remaining students are ordered by contribution to the harness.
\textsuperscript{\textdagger}Principal investigators, listed at the end in
alphabetical order by last name.

\end{thebibliography}
\end{document}
